\documentclass[11pt]{article}

\usepackage[margin=1.1in]{geometry}
\usepackage{amsmath,amssymb,amsthm,mathtools}
\usepackage{enumitem}
\usepackage{hyperref}
\usepackage[capitalise,nameinlink,noabbrev]{cleveref}

\hypersetup{colorlinks=true,linkcolor=blue,citecolor=blue,urlcolor=blue,pdftitle={Line projections do not determine complete mixability},pdfauthor={Yitzchak Shmalo}}

\newtheorem{theorem}{Theorem}[section]
\newtheorem{lemma}[theorem]{Lemma}
\newtheorem{proposition}[theorem]{Proposition}
\newtheorem{corollary}[theorem]{Corollary}
\theoremstyle{definition}
\newtheorem{definition}[theorem]{Definition}
\newtheorem{example}[theorem]{Example}
\theoremstyle{remark}
\newtheorem{remark}[theorem]{Remark}

\providecommand{\doi}[1]{\href{https://doi.org/#1}{doi:#1}}
\newcommand{\R}{\mathbb{R}}
\newcommand{\C}{\mathbb{C}}
\newcommand{\E}{\mathbb{E}}
\newcommand{\Prob}{\mathbb{P}}
\DeclareMathOperator{\supp}{supp}
\DeclareMathOperator{\re}{Re}

\title{Line projections do not determine complete mixability}
\author{Yitzchak Shmalo\\
Einstein Institute of Mathematics, The Hebrew University of Jerusalem\\
Givat Ram, Jerusalem, Israel\\
\texttt{yitzchak.shmalo@gmail.com}}
\date{}

\begin{document}
\maketitle

\begin{abstract}
A distribution $F$ on $\R^d$ is $n$-completely mixable if there are $n$ random vectors with law $F$ whose sum is
almost surely constant. Linear images of $n$-completely mixable distributions are $n$-completely mixable, and Wang
asked when the converse holds. We show that for every $n\ge 3$ there is a distribution on the plane, with a smooth
density positive on a disk or with an everywhere positive real-analytic density, whose projections onto all lines
are $n$-completely mixable while the distribution itself is not. For $n=2$ the converse holds when the mean is
finite and fails without a mean. In the other direction, one generic projection decides complete mixability for
distributions carried by compact sets of small box dimension.
\end{abstract}

\noindent\textit{2020 Mathematics Subject Classification:} 60E05, 62H05.\\
\textit{Keywords:} complete mixability, joint mixability, projections, random vectors.

\section{Introduction}

Let $n\ge 2$. A distribution $F$ on $\R^d$ is \emph{$n$-completely mixable} ($n$-CM) if there are random vectors
$X_1,\dots,X_n$, each with law $F$, such that $X_1+\dots+X_n$ is almost surely constant; the vector
$(X_1,\dots,X_n)$ is then an \emph{$n$-mix} of $F$, and when $F$ has a finite mean $b$ the constant is $nb$. The
notion arises in bounds on the risk of a sum of dependent losses. Wang and Wang \cite{WW11} characterized the
completely mixable distributions with monotone densities, Puccetti, Wang and Wang \cite{PWW12,PWW13} proved
that concave densities, and densities bounded below on an interval, are $n$-CM for $n\ge3$, and Wang and Wang
\cite{WW16} extended the theory to tuples of different distributions. Wang's survey \cite{Wang} collects the open
problems. Almost all of the theory concerns distributions on the line. R\"uschendorf and Uckelmann \cite{RU}
showed that symmetric unimodal densities are $n$-CM and made an early study of the multidimensional case, and Yin
and Zhu \cite[Theorem~4.1]{YZ} showed that a class of elliptical distributions on $\R^p$ is $n$-CM.

If $(X_1,\dots,X_n)$ is an $n$-mix of $F$ and $L\colon\R^d\to\R^k$ is linear, then $(LX_1,\dots,LX_n)$ is an
$n$-mix of the image law $F\circ L^{-1}$. This is \cite[Proposition~6]{Wang}. Wang asked for a converse
\cite[Open Problem~6]{Wang}: if $d>1$ and the projection of $F$ to every essential subspace of $\R^d$ is $n$-CM,
is $F$ $n$-CM, and if not, under which extra conditions? In the same section Wang observes that every projection of a uniform distribution on a convex
set in the plane has a concave density, so is $n$-CM for $n\ge3$ by \cite{PWW12}, and adds: ``Of course, this is not
sufficient for such a distribution to be $n$-CM on $\R^d$.'' Our first result gives distributions for which the
projections are $n$-CM and the distribution is not.

\begin{theorem}\label{thm:main}
For every $n\ge3$ there is a compactly supported distribution $P$ on $\R^2$ with mean zero such that the image of
$P$ under every linear map $\R^2\to\R$ is $n$-CM, while $P$ is not $n$-CM.
\end{theorem}

The survey does not define essential subspaces; since the proper nonzero subspaces of $\R^2$ are lines,
\cref{thm:main} answers the question in the negative for any choice of them. The failure is not caused by atoms or
by lower-dimensional supports.

\begin{theorem}\label{thm:regular}
For every $n\ge 3$ there is a distribution $P$ on $\R^2$ with mean zero such that the image of $P$ under every linear
map $\R^2\to\R$ is $n$-CM, $P$ is not $n$-CM, and $P$ has
\begin{enumerate}[label=(\alph*)]
\item a $C^\infty$ density that is positive on the interior of a closed disk and zero outside it, or
\item an everywhere positive real-analytic density on $\R^2$.
\end{enumerate}
For $n=4$, a distribution as in (a) is given explicitly in \cref{prop:disk4}.
\end{theorem}

The examples are not uniform distributions on convex sets, so they say nothing about Wang's Open Problem~5, which
asks whether such uniform distributions on $\R^d$ are $(d+1)$-CM. They are not centrally symmetric (their third or
fifth complex moments do not vanish), so they are not of the elliptically contoured or $\alpha$-symmetric kind treated in \cite{RU}.

Index two behaves differently, and here the mean matters.

\begin{proposition}\label{prop:two}
Let $X$ be a random vector in $\R^d$ with a finite mean. Then the law of $X$ is $2$-CM if and only if the law of
$\langle u,X\rangle$ is $2$-CM for every $u\in\R^d$.
\end{proposition}

\begin{example}\label{ex:stable}
There is a distribution on $\R^2$ whose image under every linear map $\R^2\to\R$ is $2$-CM but which is not $2$-CM.
\end{example}

The distribution in \cref{ex:stable} is a linear image of three independent totally skewed $1$-stable variables. It
uses the fact, studied by Puccetti, Rigo, Wang and Wang \cite{PRWW}, that without a mean the constant sum of a
mix (the \emph{center}) is no longer determined by the distribution; here the centers of the projections are
determined but are not linear in the direction.

Finally we give a positive result in the direction of Open Problem~6. It says that if the support is thin enough,
one projection suffices.

\begin{proposition}\label{prop:thin}
Let $X$ be a random vector in $\R^d$ with finite mean $b$, carried by a Borel set $S$, and let
\[
T=\{s_1+\dots+s_n-nb:\ s_1,\dots,s_n\in S\}.
\]
If $L\colon\R^d\to\R^k$ is linear and $\ker L\cap T\subseteq\{0\}$, then the law of $X$ is $n$-CM if and only if the
law of $LX$ is $n$-CM.
\end{proposition}

\begin{corollary}\label{cor:dim}
In \cref{prop:thin} let $1\le k<d$, and suppose the set of directions $E=\{t/|t|:\ 0\ne t\in T\}$ has Hausdorff
dimension less than $k$. Then for almost every $k$-dimensional subspace $V$ (with respect to the invariant
probability measure), the law of $X$ is $n$-CM if and only if its orthogonal projection onto $V$ is $n$-CM. The
hypothesis holds when $S$ is compact with upper box dimension less than $k/n$.
\end{corollary}

For example, let $K_q\subseteq[0,1]$ be the set of numbers whose base-$q$ digits are all $0$ or $1$, and let
$S=\{(x,x^2):x\in K_q\}\subseteq\R^2$. If $q>2^n$ and $F$ is a distribution carried by $S$, then for almost every direction $u$
(the exceptional set may depend on $F$), $F$ is $n$-CM if and only if its image under $x\mapsto\langle u,x\rangle$ is
$n$-CM.

\Cref{sec:four} gives the case $n=4$ of \cref{thm:main}, where the obstruction is a single moment identity, and
the explicit smooth example. \Cref{sec:three} treats $n=3$, where the scalar mixes have to be built by hand, and
\cref{sec:all} passes to all $n\ge 3$ and proves \cref{thm:regular}. \Cref{sec:two} treats index two and
\cref{sec:thin} the thin supports.

Throughout we identify $\R^2$ with $\C$, and $*$ denotes convolution. For $\phi\in\R$ write $\pi_\phi(z)=\re(e^{-i\phi}z)$, so that every
nonzero linear functional on $\R^2$ is a positive multiple of some $\pi_\phi$; since $n$-complete mixability is
preserved by scaling, it suffices to consider the maps $\pi_\phi$. We write $a(x)=1/(\pi\sqrt{1-x^2})$ for the
arcsine density on $(-1,1)$ and $T_m$ for the Chebyshev polynomial with $T_m(\cos v)=\cos(mv)$.

We use the following two standard facts.

\begin{lemma}\label{lem:basic}
\begin{enumerate}[label=(\roman*)]
\item If $F_1$ and $F_2$ are $n$-CM with the same constant sum $s$ and $0\le\lambda\le1$, then
$\lambda F_1+(1-\lambda)F_2$ is $n$-CM with sum $s$. If $X$ and $Y$ are independent random vectors whose laws are
$n$-CM, then so is the law of $X+Y$.
\item If $W$ is a random vector in $\C$ whose law is invariant under rotations about $0$, then
$(W,\omega W,\dots,\omega^{n-1}W)$ with $\omega=e^{2\pi i/n}$ is an $n$-mix of the law of $W$ with sum $0$, and so
the law of $\pi_\phi(W)$ is $n$-CM for every $\phi$.
\end{enumerate}
\end{lemma}

\begin{proof}
(i) Choose the mixture component once for the whole tuple, and in the second statement add independent mixes
coordinate by coordinate. (ii) Each $\omega^jW$ has the law of $W$, and $\sum_j\omega^j=0$.
\end{proof}

The projection of a circle law is computed as follows. If $\Theta$ has density $p(\theta)$ on $[0,2\pi)$, then
$\pi_\phi(e^{i\Theta})=\cos(\Theta-\phi)$ has density on $(-1,1)$
\begin{equation}\label{eq:proj}
\frac{p(\phi+\arccos x)+p(\phi-\arccos x)}{\sqrt{1-x^2}},
\end{equation}
the sum over the two preimages. For $p(\theta)=(1+\epsilon\cos(m\theta))/(2\pi)$ this equals
\begin{equation}\label{eq:projm}
a(x)\bigl(1+\epsilon\cos(m\phi)\,T_m(x)\bigr),
\end{equation}
because $\cos(m(\phi+v))+\cos(m(\phi-v))=2\cos(m\phi)\cos(mv)$.

\section{Four copies}\label{sec:four}

Let $\epsilon=1/10$ and let $P_4$ be the law of $Z=e^{i\Theta}$, where $\Theta$ has density
$(1+\epsilon\cos 3\theta)/(2\pi)$. Then $\E Z=0$ and $\E Z^3=\epsilon/2=1/20$.

\begin{lemma}\label{lem:cubic}
If $w_1,\dots,w_4\in\C$ and $w_1+w_2+w_3+w_4=0$, then
\[
\sum_{i=1}^4 w_i^3=3\sum_{i=1}^4\bigl(1-|w_i|^2\bigr)\prod_{j\ne i}w_j .
\]
In particular $\sum_i w_i^3=0$ if all $|w_i|=1$.
\end{lemma}

\begin{proof}
Let $e_3=\sum_i\prod_{j\ne i}w_j$ be the third elementary symmetric polynomial. By Newton's identity
$\sum_iw_i^3=e_1^3-3e_1e_2+3e_3=3e_3$. Moreover
$\sum_i|w_i|^2\prod_{j\ne i}w_j=\sum_i\overline{w_i}\prod_jw_j=\bigl(\prod_jw_j\bigr)\overline{\textstyle\sum_iw_i}=0$.
Subtracting three times this from $3e_3$ gives the identity.
\end{proof}

\begin{proof}[Proof that $P_4$ is not $4$-CM]
A $4$-mix $(Z_1,\dots,Z_4)$ of $P_4$ has sum $4\E Z=0$, and its entries lie on the unit circle. By
\cref{lem:cubic}, $\sum_iZ_i^3=0$ almost surely, so $0=4\E Z^3=1/5$, a contradiction.
\end{proof}

\begin{proposition}\label{prop:proj4}
For every $\phi$ the law of $\pi_\phi(Z)$ is $4$-CM.
\end{proposition}

\begin{proof}
By \eqref{eq:projm} the law of $Y=\pi_\phi(Z)$ has density $f(x)=a(x)(1+\epsilon cT_3(x))$ with $c=\cos3\phi$.
Since $|T_3|\le1$ on $[-1,1]$, we have $f\ge(9/10)a$. Let $h=(9/10)(a-1/4)$. It is even and nonnegative, since
$a\ge1/\pi>1/4$, and its mass is $(9/10)(1-1/2)=9/20$. A distribution with an even density is $4$-CM with sum zero:
take $(V,-V,V',-V')$ with $V,V'$ independent. The remainder $f-h$ has mass $11/20$ and mean zero, because $f$ and
$h$ both have mean zero, and
\[
\frac{f-h}{11/20}\ \ge\ \frac{(9/10)(1/4)}{11/20}=\frac{9}{22}>\frac38=\frac{3}{4\cdot 2}.
\]
By the theorem of Puccetti, Wang and Wang \cite[Theorem~3.4]{PWW13}, a distribution on an interval $[\alpha,\beta]$
with a density bounded below by $3/(n(\beta-\alpha))$ is $n$-CM for $n\ge3$. Applied with $n=4$ on $[-1,1]$, the
normalized remainder is $4$-CM with sum $0$ (the values of a density at the two endpoints can be chosen freely). Writing $f=h+(f-h)$ and using \cref{lem:basic}(i) shows that the law
of $Y$ is $4$-CM.
\end{proof}

Without removing $h$ the theorem of \cite{PWW13} does not apply to $f$ directly, since $f(0)$ is close to $1/\pi<3/8$.

To make the example smooth, for $r>0$ let $D_r$ be the rotation-invariant probability density on $\C$ proportional to
$\exp\bigl(-1/(1-|x|^2/r^2)\bigr)$ for $|x|<r$ and to $0$ otherwise. Let $\delta=1/200$, $R=1+\delta$ and
$\eta=1/1000$.

\begin{proposition}\label{prop:disk4}
Let $W$ be distributed as $Z+B$ with probability $1-\eta$, where $Z\sim P_4$ and $B$ has density $D_\delta$
independent of $Z$, and with density $D_R$ with probability $\eta$. The law of $W$ has a $C^\infty$ density, positive
on $\{|x|<R\}$ and zero outside, every image of it under a linear map $\R^2\to\R$ is $4$-CM, and it is not $4$-CM.
\end{proposition}

\begin{proof}
The density is $(1-\eta)\,P_4*D_\delta+\eta D_R$, which is smooth, supported in $\{|x|\le R\}$, and positive on the
open disk because $D_R$ is. Its projections are $4$-CM by \cref{prop:proj4} and \cref{lem:basic}: the projection of
$Z+B$ is an independent sum of $\pi_\phi(Z)$ and $\pi_\phi(B)$, both $4$-CM with sum $0$, and the projection of the
$D_R$ branch is $4$-CM with sum $0$.

Suppose $(W_1,\dots,W_4)$ is a $4$-mix. Its sum is $4\E W=0$. By \cref{lem:cubic} and $|W_i|\le R$,
\[
\Bigl|\sum_iW_i^3\Bigr|\le 3R^3\sum_i\bigl|1-|W_i|^2\bigr|,
\qquad\text{so}\qquad |\E W^3|\le 3R^3\,\E\bigl|1-|W|^2\bigr| .
\]
Since $B$ and the $D_R$ branch are rotation invariant, $\E B=\E B^2=\E B^3=0$ and the $D_R$ branch has third moment
$0$; hence $\E W^3=(1-\eta)\E Z^3=(1-\eta)/20=0.04995$. On the first branch
$\bigl||Z+B|^2-1\bigr|\le2\delta+\delta^2=401/40000$, and on the second $\bigl|1-|W|^2\bigr|\le\max(1,R^2-1)=1$.
Therefore
\[
3R^3\,\E\bigl|1-|W|^2\bigr|\le3\Bigl(\frac{201}{200}\Bigr)^3\Bigl(\frac{999}{1000}\cdot\frac{401}{40000}+\frac1{1000}\Bigr)
<0.0336<0.04995,
\]
a contradiction.
\end{proof}

\section{Three copies}\label{sec:three}

For $n=3$ we use a circle law together with an atom at the origin. Let
$\beta=1/2$, $\epsilon=1/100$, and let $P_3$ be the law of a random vector $W$ that equals $0$ with probability
$1-\beta$ and equals $e^{i\Theta}$ with probability $\beta$, where $\Theta$ has density
$(1+\epsilon\cos5\theta)/(2\pi)$. Then $\E W=0$ and $\E W^5=\beta\epsilon/2\ne0$.

\begin{proof}[Proof that $P_3$ is not $3$-CM]
A $3$-mix has sum $0$. Let $(w_1,w_2,w_3)$ be a point of $(\{0\}\cup\{|z|=1\})^3$ with $w_1+w_2+w_3=0$. If no entry
is $0$, then $|w_1+w_2|=1$ forces the angle between $w_1$ and $w_2$ to be $2\pi/3$, so the entries are $w,\omega
w,\omega^2w$ with $\omega=e^{2\pi i/3}$, and $\sum_jw_j^5=w^5(1+\omega^2+\omega)=0$. If exactly one entry is $0$ the
other two are $w,-w$, and if two are $0$ all three are. In every case $\sum_jw_j^5=0$, so a $3$-mix would give
$3\E W^5=0$.
\end{proof}

\begin{proposition}\label{prop:proj3}
For every $\phi$ the law of $\pi_\phi(W)$ is $3$-CM.
\end{proposition}

By \eqref{eq:projm} the law $\mu$ of $\pi_\phi(W)$ is $(1-\beta)\delta_0$ plus the measure with density
\[
p(y)=\beta a(y)\bigl(1+\epsilon cT_5(y)\bigr)\quad\text{on }(-1,1),\qquad c=\cos5\phi .
\]
Since $T_5$ is odd, $p(y)-p(-y)=2\beta\epsilon c\,h(y)$ for $0<y<1$, where $h=aT_5$. The idea is to produce this odd
part with triples of the form $(x,-Ux,-(1-U)x)$ and to absorb the even remainder with antipodal pairs and the atom.

\begin{lemma}\label{lem:ode}
Let $F(y)=y^{-2}\int_y^1xh(x)\,dx$ and $f=h+2F$ on $(0,1)$. Then
\[
f(y)-2\int_y^1\frac{f(x)}{x}\,dx=h(y)\quad(0<y<1),
\]
\[
f(y)=\frac{a(y)\,(32y^5-34y^3+5y)}{3},\qquad |f(y)|\le\frac{25}{3}\,y\,a(y).
\]
\end{lemma}

\begin{proof}
Differentiating, $F'(y)=-2F(y)/y-h(y)/y=-f(y)/y$, and $F(y)\to0$ as $y\to1$ since $xh$ is integrable. Hence
$F(y)=\int_y^1f(x)\,dx/x$, which is the first identity. For the closed form substitute $x=\cos v$: then
$xh(x)\,dx=-\pi^{-1}\cos v\cos5v\,dv$, and integrating gives $F(y)=(8y^3-5y)\sqrt{1-y^2}/(3\pi)$, from which the
formula for $f$ follows. For the bound, note that $\int_0^1xh(x)\,dx=\pi^{-1}\int_0^{\pi/2}\cos v\cos5v\,dv=0$, so
also $F(y)=-y^{-2}\int_0^yxh(x)\,dx$. Since $|T_5(x)|=x|16x^4-20x^2+5|\le5x$ on $[0,1]$ (the quadratic $16t^2-20t+5$
in $t=x^2$ takes the values $5$, $-5/4$ and $1$ at $t=0,5/8,1$) and $a$ is increasing on $[0,1)$, we get
$|h(y)|\le5ya(y)$ and $|F(y)|\le y^{-2}\int_0^y5x^2a(x)\,dx\le(5/3)ya(y)$. Hence $|f|\le|h|+2|F|\le(25/3)ya$.
\end{proof}

\begin{proof}[Proof of \cref{prop:proj3}]
Let $g=2\beta\epsilon c f$ on $(0,1)$, and write $g=g^+-g^-$. For $x\in(0,1)$ and $U$ uniform on $[0,1]$, the
triple $\tau_x=(x,-Ux,-(1-U)x)$ has sum $0$; its reflection $-\tau_x$ also has sum $0$. The expected counting
measure of the three entries of $\tau_x$ is $\delta_x$ plus the measure with density $2/x$ on $(-x,0)$.

Consider the measure $M$ on triples obtained by integrating $\tau_x$ against $3g^+(x)\,dx$ and $-\tau_x$ against
$3g^-(x)\,dx$, and then permuting the three entries uniformly at random. Each coordinate of $M$ has, by the previous
paragraph, the measure $\nu$ with density
\[
\nu(y)=g^+(y)+2\int_y^1\frac{g^-(x)}{x}\,dx,\qquad \nu(-y)=g^-(y)+2\int_y^1\frac{g^+(x)}{x}\,dx\qquad(0<y<1).
\]
By \cref{lem:ode}, $\nu(y)-\nu(-y)=g(y)-2\int_y^1g(x)\,dx/x=2\beta\epsilon c\,h(y)=p(y)-p(-y)$. Moreover, using
$|f(x)|/x\le(25/3)a(x)$, $\int_0^1a=1/2$ and $a\ge1/\pi$,
\[
\nu(\pm y)\le|g(y)|+2\int_0^1\frac{|g(x)|}{x}\,dx\le2\beta\epsilon\cdot\frac{25}{3}\,(a(y)+1)
\le\frac{50(1+\pi)}{3}\,\beta\epsilon\,a(y)<70\beta\epsilon\,a(y).
\]
Since $p\ge\beta(1-\epsilon)a$ and $70\epsilon<1-\epsilon$, we have $\nu\le p$ on $(-1,1)$. The residual $p-\nu$ is
therefore a nonnegative even density; let $r$ be its mass. Then $r=\beta-3\int_0^1|g|$, because $M$ and its coordinate $\nu$ both have mass $3\int_0^1|g|$.

If $r>0$ let $Y$ have density $(p-\nu)/r$, and let $M'$ be the law of the triple $(Y,-Y,0)$ with its entries
permuted uniformly, multiplied by $3r/2$. Each coordinate of $M'$ is $(r/2)\bigl(\mathcal L(Y)+\mathcal L(-Y)+\delta_0\bigr)$, which is $r$ times the law
of $Y$ plus $(r/2)\delta_0$, since $Y$ and $-Y$ have the same law. Finally let $M''$ be the point mass at $(0,0,0)$ with weight $1-\beta-r/2$, which is
nonnegative because $r\le\beta\le 2(1-\beta)$. The measure $M+M'+M''$ has total mass
$3\int|g|+3r/2+1-\beta-r/2=1$, is carried by triples with sum $0$, and each coordinate has law
$\nu+(p-\nu)+(r/2+1-\beta-r/2)\delta_0=\mu$. So $\mu$ is $3$-CM.
\end{proof}

\section{All indices from three on, and smooth densities}\label{sec:all}

\begin{proof}[Proof of \cref{thm:main}]
For $n=3$ and $n=4$ take $P_3$ and $P_4$. Let $n>3$, let $A\in\C$ with $|A|>n$, and let
$Q=(3/n)P_3+(1-3/n)\delta_A$. Every image $\pi_\phi(Q)$ is $(3/n)\pi_\phi(P_3)+(1-3/n)\delta_{\pi_\phi(A)}$; taking
a $3$-mix of $\pi_\phi(P_3)$ (with sum $0$) together with $n-3$ copies of $\pi_\phi(A)$, and permuting the $n$
entries uniformly at random, gives an $n$-mix of it.

Suppose $(V_1,\dots,V_n)$ is an $n$-mix of $Q$. Its sum is $n\E V_1=(n-3)A$. On the event that exactly $j$ entries
equal $A$, the remaining $n-j$ entries lie in the closed unit disk, so $|(n-3-j)A|\le n-j\le n$; as $|A|>n$ this
forces $j=n-3$. Hence almost surely exactly three entries differ from $A$, and they sum to $0$. Let $(U_1,U_2,U_3)$
be these three entries in uniformly random order, chosen independently of everything else. For a Borel set $G$ not
containing $A$, the expected number of entries of $(V_1,\dots,V_n)$ in $G$ is $n\,Q(G)=3P_3(G)$, so each $U_k$ has
law $P_3$, and $(U_1,U_2,U_3)$ is a $3$-mix of $P_3$, which was shown impossible. Translating $Q$ by $-(1-3/n)A$
gives a distribution with mean zero.
\end{proof}

\begin{lemma}\label{lem:limit}
Let $P$ be a distribution on $\R^2$ with mean zero which is not $n$-CM, and let $P_k\to P$ weakly, each $P_k$ with
mean zero. Then $P_k$ is not $n$-CM for all large $k$.
\end{lemma}

\begin{proof}
Otherwise, passing to a subsequence, let $\lambda_k$ be the law of an $n$-mix of $P_k$, a probability measure on
$(\R^2)^n$ carried by the closed set $H=\{(w_1,\dots,w_n):\sum_iw_i=0\}$. The marginals of $\lambda_k$ form a tight
family, so $\{\lambda_k\}$ is tight and has a weakly convergent subsequence. Its limit $\lambda$ is carried by $H$
and has all marginals equal to $P$, so it is the law of an $n$-mix of $P$.
\end{proof}

\begin{proof}[Proof of \cref{thm:regular}]
Let $n\ge 3$ and let $P$ be the distribution constructed for \cref{thm:main}, translated to mean zero; it is
supported in a disk $\{|x|\le\rho\}$. Let $R>\rho+1$.

(a) For $0<t<1$ let $P^{(t)}=(1-t)\,(P*D_t)+t\,D_R$. Its density is $C^\infty$, positive on $\{|x|<R\}$ and zero
outside. By \cref{lem:basic} its images under linear maps to $\R$ are $n$-CM: the image of the first branch is an
independent sum of an image of $P$ and an image of a rotation-invariant law, both $n$-CM with sum $0$, and the image
of the second branch is $n$-CM with sum $0$. As $t\to0$, $P^{(t)}\to P$ weakly, so by \cref{lem:limit} $P^{(t)}$ is
not $n$-CM for small $t$.

(b) For $\sigma>0$ let $P_\sigma=P*N(0,\sigma^2I)$. Its density
$x\mapsto(2\pi\sigma^2)^{-1}\int e^{-|x-z|^2/2\sigma^2}\,P(dz)$ is everywhere positive and, since $P$ has compact
support, extends to an entire function of the two coordinates; in particular it is real analytic. Its linear images
are $n$-CM by \cref{lem:basic}, and $P_\sigma\to P$ as $\sigma\to0$, so \cref{lem:limit} applies again.
\end{proof}

For $n=4$, \cref{prop:disk4} gives a choice of the parameters for which the conclusion of (a) is verified directly.

\begin{remark}
Embedding the plane in $\R^d$ gives, for every $d\ge2$ and $n\ge3$, a distribution on $\R^d$ whose images under all
linear maps $\R^d\to\R$ are $n$-CM while the distribution is not. We do not know whether, for $d\ge3$, there is such
a distribution all of whose projections onto hyperplanes are $n$-CM.
\end{remark}

\section{Two copies}\label{sec:two}

\begin{proof}[Proof of \cref{prop:two}]
If a distribution $G$ on $\R^k$ with mean $m$ is $2$-CM, a $2$-mix $(X_1,X_2)$ has $X_2=2m-X_1$, so $G$ is
symmetric about $m$; conversely, if $X$ has law $G$ and $2m-X$ has law $G$, then $(X,2m-X)$ is a $2$-mix. Let $b$
be the mean of $X$. If every $\langle u,X\rangle$ is $2$-CM, it is symmetric about $\langle u,b\rangle$, so
$\langle u,X-b\rangle$ and $\langle u,b-X\rangle$ have the same law for all $u$. By the Cram\'er--Wold theorem,
$X-b$ and $b-X$ have the same law, and the law of $X$ is $2$-CM. The other direction is \cite[Proposition~6]{Wang}.
\end{proof}

\begin{proof}[Construction for \cref{ex:stable}]
Let $S_1,S_2,S_3$ be independent with characteristic function
\[
\E e^{itS_j}=\exp\Bigl(-|t|-\frac{2i}{\pi}\,t\log|t|\Bigr),
\]
the totally skewed $1$-stable law \cite[Section~1.1]{ST}, and let $X=(S_1-S_3,\,S_2-S_3)$. For $(u,v)\ne0$ put
$\alpha=(u,v,-u-v)$, so that $uX_1+vX_2=\sum_j\alpha_jS_j$ and $\sum_j\alpha_j=0$. Using
$\alpha_jt\log|\alpha_jt|=\alpha_jt\log|t|+\alpha_jt\log|\alpha_j|$ (with $0\log0=0$), the terms in $t\log|t|$
cancel and
\[
\E e^{it(uX_1+vX_2)}=\exp\bigl(-\gamma|t|+i\kappa t\bigr),\qquad
\gamma=\sum_j|\alpha_j|>0,\quad \kappa(u,v)=-\frac2\pi\sum_j\alpha_j\log|\alpha_j| .
\]
So $uX_1+vX_2$ is a Cauchy variable with scale $\gamma$ and center $\kappa(u,v)$, and is $2$-CM through the pair
$(Y,2\kappa-Y)$. A Cauchy distribution is symmetric about exactly one point, its center, since the law of $2s-Y$
has characteristic function $\exp(-\gamma|t|+i(2s-\kappa)t)$.

If $(X^{(1)},X^{(2)})$ were a $2$-mix of $X$ with sum $C\in\R^2$, then $X$ and $C-X$ would have the same law, and so
each $uX_1+vX_2$ would be symmetric about $(uC_1+vC_2)/2$. This would make $\kappa$ linear. But $\kappa(1,0)=
\kappa(0,1)=0$ and $\kappa(1,1)=(4/\pi)\log2$.
\end{proof}

\section{Thin supports}\label{sec:thin}

\begin{proof}[Proof of \cref{prop:thin}]
One direction is \cite[Proposition~6]{Wang}. Conversely let $(Y_1,\dots,Y_n)$ be an $n$-mix of the law of $LX$; its
sum is $nLb$. Let $\kappa(y,\cdot)$ be a regular conditional distribution of $X$ given $LX=y$, and, given
$(Y_1,\dots,Y_n)$, draw $X_1,\dots,X_n$ conditionally independently with $X_i\sim\kappa(Y_i,\cdot)$. Then each
$X_i$ has the law of $X$, and almost surely $X_i\in S$ and $LX_i=Y_i$ for all $i$. Hence $t=\sum_iX_i-nb$ lies in
$T$ and $Lt=\sum_iY_i-nLb=0$, so $t=0$ almost surely.
\end{proof}

\begin{proof}[Proof of \cref{cor:dim}]
For a $k$-dimensional subspace $V$ the orthogonal projection onto $V$ has kernel $V^\perp$, and
$V^\perp\cap T\not\subseteq\{0\}$ exactly when $V^\perp$ contains a point of $E$. Let $V$ be random with the
invariant distribution. For a fixed unit vector $e$, $|P_Ve|^2$ has the beta distribution with parameters
$k/2$ and $(d-k)/2$ (by rotation invariance, $|P_Ve|^2$ has the law of $|P_Ge|^2$ for a fixed $k$-plane $G$ and a uniform unit vector $e$, that is, of $(N_1^2+\dots+N_k^2)/(N_1^2+\dots+N_d^2)$ with $N_i$ independent standard normal), so there is $C=C_{d,k}$ with
$\Prob(|P_Ve|\le r)\le Cr^k$ for $0<r<1/2$. If some unit vector $v\in V^\perp$ satisfies $|v-e|\le r$, then
$|P_Ve|=|P_V(e-v)|\le r$. Since $\dim_HE<k$, we can cover $E$ by balls with centers
$e_i\in E$ and radii $r_i<1/2$ with $\sum_ir_i^k$ as small as we like, and
$\Prob(V^\perp\cap E\ne\emptyset)\le C\sum_ir_i^k$. Hence $\Prob(V^\perp\cap E\ne\emptyset)=0$, so for almost every $V$ the hypothesis of
\cref{prop:thin} holds for the projection onto $V$.

If $S$ is compact with upper box dimension $s<k/n$, a cover of $S$ by $N$ balls of radius $\rho$ gives a cover of
$S+\dots+S$ by $N^n$ balls of radius $n\rho$, so $T$ has upper box dimension at most $ns$ and hence
$\dim_HT\le ns$. The map $t\mapsto t/|t|$ is Lipschitz on each annulus $\{2^{-m}\le|t|\le2^m\}$, so by countable
stability $\dim_HE\le ns<k$.
\end{proof}

For the example after \cref{cor:dim}: $K_q$ is covered by $2^m$ intervals of length $q^{-m}$, so its upper box
dimension is at most $\log2/\log q$, and $x\mapsto(x,x^2)$ is Lipschitz, so $S$ has upper box dimension at most
$\log2/\log q<1/n$ when $q>2^n$. \Cref{cor:dim} applies with $d=2$, $k=1$.

\section*{Acknowledgments}

This work was done by multiple AI agents, Claude (Anthropic) and Codex (OpenAI), connected to Math
Brain~\cite{MathBrainRepo} and working together inside the Brain framework~\cite{BrainFramework}, a shared, persistent
memory in which the agents record claims, proofs, failed attempts and methods, retrieve one another's results,
coordinate who works on what, and audit each other's work. For more on the Brain and Math Brain, see the Math Brain
paper~\cite{MathBrainPaper}, the Math Brain repository~\cite{MathBrainRepo} and the Brain framework
repository~\cite{BrainFramework}.

\section*{Funding}

The research presented in this paper was supported by the European Research Council (ERC) under the European
Union's Horizon Europe research and innovation programme (grant agreement No.\ 101041711). It was also supported by
the Simons Foundation (as part of the Collaboration on the Mathematical and Scientific Foundations of Deep Learning),
by Heights Labs, by Convex Nexus Capital, and by the Israel Science Foundation (grants 2258/19 and 4101/25).

\section*{Declaration of competing interest}

The author holds equity in Heights Labs and Convex Nexus Capital and has a management role at Convex Nexus Capital.

\section*{Declaration of generative AI and AI-assisted technologies}

The arguments and programs were developed by AI agents, Claude (Anthropic) and Codex (OpenAI), connected to Math Brain
and working under the author's direction; the agents also drafted the text. The question is Open Problem~6 of Wang
\cite{Wang}. The lower-density criterion used in \cref{sec:four} is due to Puccetti, Wang and Wang \cite{PWW13}, and
the characteristic function of the $1$-stable law is standard \cite{ST}. The author directed the work and is
responsible for the content of the paper and for decisions concerning its dissemination.

\section*{Data and code availability}

The program \texttt{check\_constants.py} in the repository
\begin{center}\url{https://github.com/yspennstate/projection-complete-mixability/tree/v1.0}\end{center}
recomputes the constants of \cref{prop:proj4,prop:disk4} in exact rational arithmetic and checks
\cref{lem:ode} and the inequality $\nu\le p$ of \cref{sec:three} numerically on a grid.

\begin{thebibliography}{99}

\bibitem{PRWW}
G.~Puccetti, P.~Rigo, B.~Wang and R.~Wang, Centers of probability measures without the mean, \emph{J. Theoret.
Probab.} 32 (2019), 1482--1501.

\bibitem{PWW12}
G.~Puccetti, B.~Wang and R.~Wang, Advances in complete mixability, \emph{J. Appl. Probab.} 49 (2012), 430--440.

\bibitem{PWW13}
G.~Puccetti, B.~Wang and R.~Wang, Complete mixability and asymptotic equivalence of worst-possible VaR and ES
estimates, \emph{Insurance Math. Econom.} 53 (2013), 821--828.

\bibitem{RU}
L.~R\"uschendorf and L.~Uckelmann, Variance minimization and random variables with constant sum, in
\emph{Distributions with Given Marginals and Statistical Modelling}, Kluwer, Dordrecht, 2002, 211--222.

\bibitem{ST}
G.~Samorodnitsky and M.~S.~Taqqu, \emph{Stable Non-Gaussian Random Processes}, Chapman and Hall, New York, 1994.

\bibitem{Wang}
R.~Wang, Current open questions in complete mixability, \emph{Probab. Surv.} 12 (2015), 13--32.
\doi{10.1214/14-PS250}

\bibitem{WW11}
B.~Wang and R.~Wang, The complete mixability and convex minimization problems with monotone marginal densities,
\emph{J. Multivariate Anal.} 102 (2011), 1344--1360.

\bibitem{WW16}
B.~Wang and R.~Wang, Joint mixability, \emph{Math. Oper. Res.} 41 (2016), 808--826.

\bibitem{YZ}
C.~Yin and D.~Zhu, Joint mixability of elliptical distributions and related families, preprint (2023),
arXiv:1706.05499v4.

\bibitem{MathBrainRepo}
Y.~Shmalo, \emph{Math Brain: public snapshot, code and tests}, 2026.
\url{https://github.com/yspennstate/math-brain-public}

\bibitem{MathBrainPaper}
Y.~Shmalo, \emph{Math Brain: shared, auditable memory for multi-agent mathematical research}, October 2026.
\url{https://github.com/yspennstate/math-brain-public/tree/main/papers/math-brain-2026-10-24}

\bibitem{BrainFramework}
Y.~Shmalo, \emph{Brain Framework}, public software repository, 2026.
\url{https://github.com/yspennstate/brain-framework-public}

\end{thebibliography}

\end{document}
